\chapterHeading{Data structures}

\sectionHeading{Segment tree (point update, range query)}
\begin{lstlisting}[language=C++]
#include <iostream>
#include <vector>
 
// build            (1, 1, n    );
// update           (1, 1, n,   l, r,   {a, b});
// Info ans = query (1, 1, n,   l, r);
// make sure to convert into 0 based indexing
 
struct Info {
    int sum=0;
};
 
const int MAXN = 2e5+5;
 
vector<Info> segtree(4 * MAXN);
vector<int> v(MAXN);
pair<int, int> lazy[4 * MAXN];  // pending operations at the node (**for children**)
 
Info Identity(){
    Info a;
    return a;
}
 
pair<int, int> updateIdentity() {
    return {1, 0};
}
 
Info merge(Info &a, Info &b){
    Info res;
    res.sum = a.sum + b.sum;
    return res;
}
void applyUpdate(int id, int l, int r, pair<int, int>lazy_val) {
    // lazy_val.first = a, lazy_val.second = b
    // Every element x becomes a*x + b
    // Sum becomes a*sum + b*(r-l+1)
 
    segtree[id].sum = lazy_val.first*segtree[id].sum + lazy_val.second*(r-l+1);
 
    auto temp = lazy[id];
 
    /// lazy[id] = (a, b) -> ax + b
    /// v = (p, q)       -> px + q
    /// New operation: p(ax + b) + q
    ///                 = (pa)x + (pb + q)
 
    lazy[id].first = lazy_val.first * temp.first;
    lazy[id].second = lazy_val.first * temp.second + lazy_val.second;
}
 
void lazypush(int id, int l, int r) {
    // no change in tree[id]
    if (lazy[id] == updateIdentity()) return;
 
    if (l == r) {
        lazy[id] = updateIdentity();
        return;
    }
 
    int lc = 2 * id;
    int rc = 2 * id + 1;
    int m = (l + r) / 2;
 
    applyUpdate(lc, l, m, lazy[id]);
    applyUpdate(rc, m+1, r, lazy[id]);
 
    lazy[id] = updateIdentity();
}
 
// O(n) // lazy build
void build(int id, int l, int r) {
    lazy[id] = updateIdentity();
 
    if (l == r) {
        segtree[id].sum = v[l];
        return;
    }
 
    int lc = 2 * id;
    int rc = 2 * id + 1;
    int m = (l + r) / 2;
 
    build(lc, l, m);
    build(rc, m+1, r);
 
    segtree[id] = merge(segtree[lc], segtree[rc]);
}
 
// O(log n)
Info query(int id, int l, int r, int L, int R) {
    lazypush(id, l, r);
 
    if (L <= l && r <= R) {
        // perfect -> [l, r] is inside [L, R]
        return segtree[id];
    }
 
    if (r < L || R < l) {
        // outside -> [l, r] does not overlap with [L, R]
        return Identity();
    }
 
    // extra -> remaining
    int lc = 2 * id;
    int rc = 2 * id + 1;
    int m = (l + r) / 2;
 
    auto ans1 = query(lc, l, m, L, R);
    auto ans2 = query(rc, m+1, r, L, R);
 
    return merge(ans1, ans2);
}
 
 
 
// O(log n)
void update(int id, int l, int r, int L, int R, pair<int, int> v) {
    lazypush(id, l, r);
 
    if (L <= l && r <= R) {
        // perfect -> [l, r] is inside [L, R]
        applyUpdate(id, l, r, v);
        return;
    }
 
    if (r < L || R < l) {
        // outside -> [l, r] does not overlap with [L, R]
        return;
    }
 
    int lc = 2 * id;
    int rc = 2 * id + 1;
    int m = (l + r) / 2;
 
    update(lc, l, m, L, R, v);
    update(rc, m+1, r, L, R, v);
 
    segtree[id] = merge(segtree[lc], segtree[rc]);
}
 
void debug(int id, int l, int r){
    int mid = (l+r)/2;
    int lc = 2*id;
    int rc = 2*id + 1;
 
    cout << id << " -- [" << l << ", " << r << "] --> " << segtree[id].sum << nl;
 
    if(l == r) return;
 
    debug(lc, l, mid);
    debug(rc, mid+1, r);
}
 
void solve(){
    int n, q; cin>>n>>q;
 
    for(int i=1; i<=n; i++) cin>>v[i];
 
    build(1, 1, n);
 
    while(q--){
        int k; cin>>k;
 
        if(k==1){
            // Increase each value in [a, b] by x
            int a, b, x; cin>>a>>b>>x;
            update(1, 1, n, a, b, {1, x});
        }
        else if(k==2){
            // Set each value in [a, b] to x
            int a, b, x; cin>>a>>b>>x;
            update(1, 1, n, a, b, {0, x});
        }
        else{
            // Calculate the sum of values in [a, b]
            int a, b; cin>>a>>b;
            cout << query(1, 1, n, a, b).sum << nl;
        }
    }
}

int32_t main(){
    fast

    int t = 1, w = 1;
    // cin >> t;

    while(t--){
        // cout << "Case " << w++ << ": ";
        solve();
    }
}
\end{lstlisting}

\sectionHeading{MO's algorithm}
\begin{lstlisting}[language=C++]
struct tup{
    int block, r, l, i;
};
int n, q, distinct;
vector<tup>queries;
vector<int>cs, ans, freq;
void add (int pos){
    int x = cs[pos];
    if (freq[x] == 0) ++distinct;
    ++freq[x];
}
void remove (int pos){
    int x = cs[pos];
    --freq[x];
    if (freq[x] == 0) --distinct;
}
void solve(){
    cin >> n >> q;
    queries.resize(q); cs.resize(n); ans.resize(q); freq.resize(n);
    for (int i = 0; i < n; ++i) cin>>cs[i];
 
    vector <int>coo = cs;
    sort(coo.begin(), coo.end());
    coo.erase(unique(coo.begin(), coo.end()), coo.end());

    for(int i=0; i<n; i++)
        cs[i] = lower_bound(coo.begin(), coo.end(), cs[i]) - coo.begin();
 
    for(int i=0; i<q; i++){
        int l, r; cin>>l>>r;
        --l, --r;
        int block = l/sqrt(n);
        queries[i] = {block, (block & 1) ? -r : r, l, i};
    }
 
    sort(all(queries), [](tup &a, tup &b){
        if(a.block != b.block) return a.block < b.block;
        if(a.r != b.r) return a.r < b.r;
        return a.l < b.l;
    });

    int L=0, R=-1;
    for(int it=0; it<q; it++){
        auto [block, r, l, i] = queries[it];
        r = abs(r);
        while(R < r) add(++R);
        while(L > l) add(--L);
        while(L < l) remove(L++);
        while(R > r) remove(R--);
 
        ans[i] = distinct;
    }
 
    for(int i=0; i<q; i++) cout<<ans[i]<<nl;
}

\end{lstlisting}


